% ============================================================
%  Razões Trigonométricas — Apostila Premium | PratiqueJá
%  Engine: XeLaTeX
% ============================================================
\documentclass[12pt]{article}
\usepackage{../../pratiqueja}
\usepackage{tikz}
\usetikzlibrary{calc}

% \hlm — destaque de resultado em modo-math (cor sem bold)
\newcommand{\hlm}[1]{{\color{darkblue}#1}}

\tikzset{
  tri/.style={draw=darkblue, line width=1.2pt, fill=accentblue!12!white},
  trj/.style={draw=darkblue, line width=1.2pt, fill=badgegreen!12!white},
  ra/.style={draw=darkblue, line width=0.8pt},
  cl/.style={font=\footnotesize\bfseries},
}

\setsubject{Razões Trigonométricas}

\begin{document}
\pagenumbering{arabic}
\setstretch{1.2}
\color{bodytext}

\heroheader{Razões Trigonométricas}%
           {Seno, Cosseno, Tangente e Ângulos Notáveis}%
           {Matemática · Intermediário}

\setlength{\columnsep}{18pt}
\setlength{\columnseprule}{0.5pt}
\def\columnseprulecolor{\color{exborder}}

\begin{multicols}{2}

%─────────────────────────────────────────────────────────────
\TW{Def}{badgepurple}{Razões Trigonométricas}{Relação entre lados e ângulos do triângulo retângulo}

\regra{As razões trigonométricas relacionam os \textbf{lados} de um
triângulo retângulo aos seus \textbf{ângulos agudos}.}

\begin{center}\vspace{2pt}
\begin{tikzpicture}[scale=0.85]
\coordinate (A) at (0,0); \coordinate (B) at (4,0); \coordinate (C) at (0,2.8);
\draw[tri] (A)--(B)--(C)--cycle;
\draw[ra] (0.4,0)--(0.4,0.4)--(0,0.4);   % ângulo reto em A
\node[cl] at (2,-0.34){$ca$};
\node[cl] at (-0.42,1.4){$co$};
\node[cl] at (2.4,1.6){$h$};
% ângulo alpha em B
\draw[graycolor, line width=0.7pt] ([shift=(180:0.62)]B) arc (180:145:0.62);
\node[font=\scriptsize] at ($(B)+(162:0.98)$){$\alpha$};
\end{tikzpicture}
\end{center}

\obs{\textbf{Hipotenusa} $(h)$: lado oposto ao ângulo reto — o maior lado.\\[2pt]
\textbf{Cateto oposto} $(co)$: lado oposto ao ângulo $\alpha$.\\[2pt]
\textbf{Cateto adjacente} $(ca)$: lado que forma o ângulo $\alpha$ junto
com a hipotenusa.}

%─────────────────────────────────────────────────────────────
\TW{f($\alpha$)}{darkblue}{Definição das Razões}{Seno, cosseno e tangente}

\formula{$\operatorname{sen}\alpha = \dfrac{co}{h} \qquad
\cos\alpha = \dfrac{ca}{h}$\\[2pt]
$\tan\alpha = \dfrac{co}{ca}$}

\nota{\textbf{Macete SOH-CAH-TOA:}\\
\textbf{S}eno = \textbf{O}posto / \textbf{H}ipotenusa\\
\textbf{C}osseno = \textbf{A}djacente / \textbf{H}ipotenusa\\
\textbf{T}angente = \textbf{O}posto / \textbf{A}djacente}

\SH{Relação entre as razões}
\exemp{%
  $\tan\alpha = \dfrac{\operatorname{sen}\alpha}{\cos\alpha}$ \quad
  {\footnotesize a tangente é o quociente entre seno e cosseno.}%
}

%─────────────────────────────────────────────────────────────
\TW{ID}{badgepurple}{Identidades Fundamentais}{Relações sempre verdadeiras}

\SH{Identidade Pitagórica}
\formula{$\operatorname{sen}^2\alpha + \cos^2\alpha = 1$}
\obs{Decorre de Pitágoras $co^2 + ca^2 = h^2$. Dividindo por $h^2$:
$\left(\dfrac{co}{h}\right)^2 + \left(\dfrac{ca}{h}\right)^2 = 1$.}

\SH{Ângulos Complementares}
\obs{Em um triângulo retângulo, os dois ângulos agudos somam $90°$
(são complementares). Logo:}
\exemp{%
  $\operatorname{sen}\alpha = \cos(90°-\alpha)$\\[2pt]
  $\cos\alpha = \operatorname{sen}(90°-\alpha)$\\[2pt]
  $\tan\alpha = \dfrac{1}{\tan(90°-\alpha)}$%
}
\nota{$\operatorname{sen}30° = \cos60° = \dfrac{1}{2}$ \quad e \quad
$\cos30° = \operatorname{sen}60° = \dfrac{\sqrt{3}}{2}$.}

%─────────────────────────────────────────────────────────────
\TW{Ex}{badgegreen}{Exemplo}{Triângulo com lados 6, 8 e 10}

\begin{center}\vspace{2pt}
\begin{tikzpicture}[scale=0.85]
\coordinate (A) at (0,0); \coordinate (B) at (4,0); \coordinate (C) at (0,3);
\draw[trj] (A)--(B)--(C)--cycle;
\draw[ra] (0.4,0)--(0.4,0.4)--(0,0.4);
\node[cl] at (2,-0.34){$8$};
\node[cl] at (-0.34,1.5){$6$};
\node[cl] at (2.45,1.7){$10$};
\draw[graycolor, line width=0.7pt] ([shift=(180:0.62)]B) arc (180:143:0.62);
\node[font=\scriptsize] at ($(B)+(161:0.98)$){$\alpha$};
\end{tikzpicture}
\end{center}

\exemp{%
  $\operatorname{sen}\alpha = \dfrac{6}{10} = \hlm{\dfrac{3}{5}}$ \qquad
  $\cos\alpha = \dfrac{8}{10} = \hlm{\dfrac{4}{5}}$\\[3pt]
  $\tan\alpha = \dfrac{6}{8} = \hlm{\dfrac{3}{4}}$%
}
\nota{\textbf{Verificação:}
$\left(\dfrac{3}{5}\right)^2 + \left(\dfrac{4}{5}\right)^2
= \dfrac{9}{25} + \dfrac{16}{25} = \dfrac{25}{25} = 1$ \ok}

%─────────────────────────────────────────────────────────────
\TW{Tab}{darkblue}{Ângulos Notáveis}{$30°$, $45°$ e $60°$}

\obs{Os valores mais cobrados em provas. O seno segue a sequência
$\dfrac12,\ \dfrac{\sqrt2}{2},\ \dfrac{\sqrt3}{2}$ e o cosseno é o
complemento (ordem inversa).}

\begin{center}\vspace{3pt}
\setlength{\tabcolsep}{12pt}\renewcommand{\arraystretch}{1.6}\small
\begin{tabular}{c|ccc}
\rowcolor{rulebg}
$\alpha$ & $30°$ & $45°$ & $60°$\\
\hline
\textbf{sen} & $\dfrac12$ & $\dfrac{\sqrt2}{2}$ & $\dfrac{\sqrt3}{2}$\\
\rowcolor{exbg}
\textbf{cos} & $\dfrac{\sqrt3}{2}$ & $\dfrac{\sqrt2}{2}$ & $\dfrac12$\\
\textbf{tan} & $\dfrac{\sqrt3}{3}$ & $1$ & $\sqrt3$\\
\end{tabular}
\end{center}

\nota{\textbf{Macete:} para o seno, escreva
$\dfrac{\sqrt0}{2},\ \dfrac{\sqrt1}{2},\ \dfrac{\sqrt2}{2},\ \dfrac{\sqrt3}{2},\ \dfrac{\sqrt4}{2}$
para $0°, 30°, 45°, 60°, 90°$. Para o cosseno, leia ao contrário.}

%─────────────────────────────────────────────────────────────
\TW{Ex+}{badgegreen}{Problemas Contextualizados}{Aplicando ângulos notáveis}

\SH{P1 — cateto oposto ($\operatorname{sen}30°$)}
\begin{center}\vspace{1pt}
\begin{tikzpicture}[scale=0.85]
\coordinate (A) at (0,0); \coordinate (B) at (4,0); \coordinate (C) at (0,2.31);
\draw[tri] (A)--(B)--(C)--cycle;
\draw[ra] (0.38,0)--(0.38,0.38)--(0,0.38);
\node[cl] at (-0.4,1.2){$x$};
\node[cl] at (2.45,1.45){$16$};
\draw[graycolor, line width=0.7pt] ([shift=(180:0.7)]B) arc (180:150:0.7);
\node[font=\scriptsize] at ($(B)+(165:1.05)$){$30°$};
\end{tikzpicture}
\end{center}
\exemp{%
  Hipotenusa $=16$, achar o cateto oposto $x$. Como
  $\operatorname{sen}30° = \dfrac12$:\\[2pt]
  $\dfrac12 = \dfrac{x}{16} \Rightarrow 2x = 16 \Rightarrow x = \hlm{8}$%
}

\SH{P2 — hipotenusa ($\cos60°$)}
\exemp{%
  Ângulo $60°$, cateto adjacente $=5$, achar a hipotenusa $h$.
  Como $\cos60° = \dfrac12$:\\[2pt]
  $\dfrac12 = \dfrac{5}{h} \Rightarrow h = 5\cdot2 = \hlm{10}$%
}

\SH{P3 — cateto adjacente ($\tan45°$)}
\exemp{%
  Ângulo $45°$, cateto oposto $=7$, achar o cateto adjacente $ca$.
  Como $\tan45° = 1$:\\[2pt]
  $1 = \dfrac{7}{ca} \Rightarrow ca = \hlm{7}$%
}
\nota{Quando $\tan\alpha = 1$, o triângulo é \textbf{isósceles} (catetos
iguais) — ocorre exatamente em $45°$.}

\end{multicols}

\end{document}
